% ECONOMICS_PROBLEM: {"id": "J24-585", "title": "Retraining after occupational displacement", "jel_code": "J24", "source_id": "OP-0353"}
% !TeX program = xelatex
\documentclass[11pt,a4paper]{article}
\usepackage[margin=23mm]{geometry}
\usepackage{fontspec}
\setmainfont{Latin Modern Roman}
\usepackage{amsmath,amssymb,mathtools}
\usepackage{xcolor,enumitem,booktabs,longtable,array,xurl,hyperref}
\definecolor{muted}{HTML}{53636C}
\hypersetup{colorlinks=true,urlcolor=blue,linkcolor=blue}
\setlength{\parindent}{0pt}
\setlength{\parskip}{5pt}
\newcommand{\R}{\mathbb R}
\newcommand{\E}{\mathbb E}
\newcommand{\N}{\mathbb N}
\newcommand{\ind}{\mathbf 1}
\newcommand{\Var}{\operatorname{Var}}
\newcommand{\Cov}{\operatorname{Cov}}
\newcommand{\supp}{\operatorname{supp}}
\DeclareMathOperator*{\argmin}{arg\,min}
\DeclareMathOperator*{\argmax}{arg\,max}
\newcommand{\entrymeta}[3]{{\sffamily\small\color{muted}\textbf{\texttt{#1}}\quad|\quad #2\par #3\par}}
\newcommand{\Problemref}[1]{\texttt{#1}}
\newcommand{\JELref}[2]{\texttt{#2} (JEL \texttt{#1})}
\begin{document}
\section*{Retraining after occupational displacement}
\textbf{Problem ID:} \texttt{J24-585}\quad \textbf{JEL:} \texttt{J24}\par
% BEGIN ECONOMICS STATEMENT
\label{problem:OP-0353}
{\sffamily\small\color{muted}Original subject group: Labor markets and work\par}
\label{definitions:problem:30007635}
\textbf{Audit classification:} Published research agenda. Metadata and full staff report checked. Workplace training and other occupational populations are future directions; results use matched comparisons, not randomized training.
\subsection*{Definitions and precise research target}
\textbf{Complete editorial problem.} Fix a population of workers displaced from shrinking occupations and a finite menu \(\mathcal P\) of retraining programs. Program \(p\) specifies curriculum, duration, financial support, employer links, and delivery mode; \(p=0\) denotes the declared job-search assistance comparator. Let \(X_i\) contain worker \(i\)'s age, prior skills, occupation, and local demand at enrollment. Let \(Y_{it}(p)\) be real earnings at date \(t\), including zero earnings during nonemployment. With discount factor \(\delta\in(0,1]\), define \(V_i(p)=\sum_{t=0}^{T}\delta^tY_{it}(p)\), where \(T\) is a preregistered follow-up horizon. Estimate \(\tau_p(x)=E[V_i(p)-V_i(0)\mid X_i=x]\), program costs, and stable-employment outcomes. Determine a feasible assignment rule \(\pi_p(X_i)\ge0\), \(\sum_{p\in\mathcal P}\pi_p(X_i)=1\) maximizing expected earnings gains subject to a fixed training budget. Distinguish skill acquisition from credential signals and occupation selection through additional experimental variation. State whether randomization identifies participation or offer effects, and address compliance and differential attrition. Transportability to white-collar workers, workplace training, changing technology, and weaker labor demand requires evidence rather than an assumption that a favorable short-run estimate persists.

\textbf{Definitions and admissible scope.} Programs are assignments or offers, so \(p=0\) and every active program have specified access, uptake, and completion rules. Fix \(T<\infty\), \(\delta\in(0,1]\), integrable earnings including zero in unemployment, and per-person resource cost \(c_p(X)\ge0\). A randomized assignment rule consists of measurable probabilities \(\pi_p(x)\ge0\) summing to one. Its earnings value is \(\sum_pE[\pi_p(X)\tau_p(X)]\), with \(\tau_0=0\), and feasibility means \(\sum_pE[\pi_p(X)c_p(X)]\le B\). Employment stability means, for example, paid employment in each of the last \(1\le s\le T+1\) periods; fix \(s\) in advance. Offer randomization identifies intention-to-treat effects under consistency and follow-up; effects of completion require additional exclusion and compliance restrictions. Unknown treatment effects and transportability, rather than existence of a maximum in a known finite menu, are the research task. The optimization is explicitly over randomized rules \(\pi=(\pi_p)_{p\in\mathcal P}\), where \(\pi_p:\mathcal X\to[0,1]\) is measurable, \(\sum_p\pi_p=1\), and \(\pi_0\) is comparator assignment probability. Its value is \(\sum_pE[\pi_p(X)\tau_p(X)]\), and its expected resource constraint is \(\sum_pE[\pi_p(X)c_p(X)]\le B\); the deterministic rule is the special case with probabilities zero or one. A stable-employment outcome is \(\prod_{t=T-s+1}^T1\{h_{it}(p)>0\}\), where \(1\le s\le T+1\) is fixed and \(h\) denotes paid hours. Conditional effects are asserted only almost everywhere on common assignment support. Under observational enrollment, the stipulated selection restrictions, rather than the source study alone, determine identification.
\subsection*{Variants and assumption changes}
\begin{enumerate}
\item Workplace and occupation version (source-motivated): compare on-the-job retraining with classroom retraining for white-collar workers, specifying employer subsidy and retention rules. This changes both the population and the treatment.
\item Demand-state version (source-motivated extension): permit \(\tau_p(x,s)\) and costs to depend on labor-demand state \(s\). Identify state-specific assignment policies and test whether skill and credential effects persist when vacancies are scarce.
\end{enumerate}
\subsection*{References and source audit}
\textbf{Support and access.} Metadata and full staff report checked. Workplace training and other occupational populations are future directions; results use matched comparisons, not randomized training. \textbf{Locator:} Section 6, Discussion, pp. 15--16, last paragraph.
\begin{enumerate}\raggedright
\item\hypertarget{definitions:source:30007635:1}{} Benjamin Hyman; Benjamin Lahey; Karen Ni; Laura Pilossoph. 2025. \emph{How Retrainable Are AI-Exposed Workers?}. Conclusion, printed pp. 15–16, final paragraph.
\url{https://www.newyorkfed.org/medialibrary/media/research/staff_reports/sr1165.pdf}.
DOI: \url{https://doi.org/10.59576/sr.1165}.
\end{enumerate}

\subsection*{Common conventions: Source mathematical conventions}

These conventions apply to each entry unless explicitly replaced.

\textbf{Models and identification.} A primitive model \(m\) specifies all probability laws, feasible choices, information, equilibrium and measurement rules used by its target. The admissible class \(\mathcal M\) is fixed independently of outcomes. If \(P_O(m)\) is its observable probability law and \(\tau(m)\) the target, its identified set at observable law \(P\) is
\[
 \mathcal I_\tau(P)=\{\tau(m):m\in\mathcal M,\ P_O(m)=P\}.
\]
Point identification means that this set is a singleton. An empty set signals incompatibility of data and assumptions. Coordinate bounds need not describe the joint set. Bounds are sharp only if every retained target is attainable or arbitrarily approachable by compatible models. Finite bounds require explicit restrictions on unobserved quantities; unrestricted confounding, hidden wealth, extrapolation or arbitrary equilibrium selection can yield uninformative sets.

\textbf{Causal interventions.} A potential outcome \(Y_i(p)\) is the outcome under a complete policy regime \(p\), including timing, financing, information and market spillovers. The target population and units are fixed before treatment. With interference, use the complete assignment vector or a justified exposure mapping; individual no-interference notation alone cannot identify an economy-wide intervention. A difference of mean potential outcomes does not require the cross-world joint distribution. Conditioning on post-treatment migration, survival, enrollment or employment changes the estimand and needs separate justification.

\textbf{Statistical statements.} An estimator, confidence set, forecast or test specifies a sampling law, model class and loss/coverage criterion. Uniform \((1-\alpha)\)-coverage means \(\inf_{m\in\mathcal M}\Pr_m\{\tau(m)\in C_n\}\geq1-\alpha\), or an explicitly asymptotic version. The whole parameter space trivially has coverage, so informativeness requires a stated width, risk or power objective. A forecasting improvement is relative to a named benchmark, information set and evaluation population.

\textbf{Equilibrium and optimization.} An equilibrium comprises feasible plans, optimality, beliefs where needed, budget constraints and all market/resource clearing. The primitives or a stated benchmark define each correspondence \(\mathcal E(p;m)\). Nonemptiness is assumed or proved, never inferred just from compact policy sets. A selected-equilibrium target uses a declared selection rule; a robust target quantifies over all permitted equilibria. Use suprema or infima unless attainment is established. An \(\varepsilon\)-optimal policy is feasible and within \(\varepsilon\) of the finite optimum value. Report an empty feasible set separately.

\textbf{Welfare and accounting.} Normative utility, interpersonal normalization, social weights, numeraire, discounting and the use of public receipts are primitives. Behavioral utility need not coincide with the welfare criterion. Household welfare includes ownership income and rents once; adding profits again to compensating variation generally double-counts them. A monetary equivalent is defined by an explicit transfer and price convention, with monotonicity and range conditions. Average effects, mechanism contributions, transfers and welfare losses are different targets. Infinite sums require convergence and dynamic feasible plans require terminal or no-Ponzi conditions.

\textbf{Source versions.} A working-paper date, revision date and journal publication year are distinguished. A DOI for a journal version is not automatically the DOI of an earlier manuscript. Locators refer to the stated version and printed pages where specified. A metadata-only check does not verify a claim about a full-text conclusion. Every entry's source subsection narrows the provenance claim accordingly.
% END ECONOMICS STATEMENT
\end{document}
